\section{CREMI Sample A Results}

The CREMI Drosophila EM dataset serves as a correctness benchmark. Because the
input is human-annotated labels rather than automated segmentation output, the
ground-truth set contains only genuine split errors (none of the ambiguity introduced by
an imperfect automated segmenter). This makes CREMI well-suited for verifying
that the validation rules reject false positives while accepting true positives.

Table~\ref{tab:cremi} summarizes the final CREMI results after parameter tuning
(CompositeScoreRule threshold = 0.15, corresponding to Experiments 5--7 in the
development log).

\begin{table}[H]
\centering
\caption{CREMI Sample A final results (64$\times$256$\times$256 voxels, 40$\times$4$\times$4 nm resolution).}
\label{tab:cremi}
\begin{tabular}{@{}lrr@{}}
\toprule
Metric & Value & Notes \\
\midrule
Fragments & 243 & After stitching \\
Candidates generated & 3{,}369 & \\
Accepted & 20 & \\
Rejected & 3{,}339 & \\
Ambiguous & 10 & \\
True Positives & 20 & \\
False Positives & 0 & \\
False Negatives & 2 & \\
Precision & 1.000 & \\
Recall & 0.909 & 20 / 22 \\
F1 & 0.952 & \\
\bottomrule
\end{tabular}
\end{table}

The two false negatives are pairs whose composite scores fall below the
\texttt{CompositeScoreRule} threshold of 0.65 --- one short-gap pair with
alignment-and-continuity-default scoring (composite $\approx 0.38$), and one
borderline pair whose ground-truth neuron bends sharply enough that the
CurvatureRule correctly flags it as implausible. The perfect precision (zero
false positives) confirms that the conservative rule set does not over-merge.

\section{XPRESS Training Volume: Experiment Progression}

The XPRESS challenge is the primary domain-appropriate benchmark. Experiments
8--24 trace a systematic development process starting from a baseline of
Coverage = 23\%, Recall = 0.160, and ending at Coverage = 78.7\%,
Recall = 0.9958. Table~\ref{tab:xpress_progress} summarizes key experiments.

\begin{table}[H]
\centering
\caption{XPRESS training volume experiment progression. Each row shows the key
change from the previous experiment, ground-truth coverage, and recall on covered pairs.}
\label{tab:xpress_progress}
\small
\begin{tabularx}{\textwidth}{@{}l>{\RaggedRight}Xrr>{\RaggedRight}X@{}}
\toprule
Exp & Key change & Cov. & Recall & Notes \\
\midrule
8  & Endpoint graph, 300$^3$ crop & 23.0\% & 0.160 & Interior splits invisible \\
9  & Skeleton-node graph, 699$^3$ & 79.8\% & 0.605 & +56.8 pp coverage \\
10 & \texttt{max\_skeleton\_voxels} = 500K & 62.1\% & 0.851 & Better skeletons \\
11 & Long-range endpoint pass (2000 nm) & 81.6\% & 0.798 & \\
12 & Threshold tuning, proximity cliff removed & 81.7\% & 0.898 & \\
13 & 3000 nm radius & --- & 0.865 & \textit{Negative; reverted} \\
14 & Distance-conditioned weights & 83.5\% & 0.889 & Long-range fix \\
15 & Long-range accept threshold & 83.5\% & 0.889 & No effect \\
16 & CurvatureRule skip for long-range & 83.5\% & 0.916 & \\
17 & MaxDistance 20K nm, SizeDiscrep 15.0 & 83.5\% & 0.994 & Prior best \\
\midrule
18 & Partner limit (AND semantics) & 83.5\% & 0.008 & \textit{Negative; reverted} \\
19 & MinGapRule (150 nm) & 75.9\% & 0.966 & \textit{Negative on recall} \\
20 & ML filter (gradient boosting) & 83.5\% & 0.852 & Precision 0.037 at recall 0.85 \\
21 & 3000 nm + ML filter & --- & --- & \textit{Negative; ML fails at recall $\geq$ 0.99} \\
22 & FN diagnosis (no code change) & --- & --- & Diagnostic only \\
\midrule
23 & Held-out validation (Exp 17 config) & --- & --- & Validation: 72.4\% cov., Recall=0.885 \\
24 & Direction weights + MinGapRule off & \textbf{78.7\%} & \textbf{0.9958} & Best; FN=5 (all CurvatureRule) \\
25 & CurvatureRule skip 1000$\to$1100 nm & 78.7\% & 0.9949 & \textit{Negative; reverted} \\
\bottomrule
\end{tabularx}
\end{table}

\paragraph{Key turning points.}
Three changes account for most of the recall improvement:

\begin{enumerate}
  \item \textbf{Experiment 9 (skeleton-node graph):} The switch from endpoint-only
        to full skeleton-node indexing raised coverage from 23\% to 79.8\% in a
        single change. The endpoint graph is structurally incapable of detecting
        splits that occur in the interior of long axons; the skeleton-node graph
        can detect any split regardless of its position along the axon.

  \item \textbf{Experiment 14 (distance-conditioned scoring):} Proximity score
        decays exponentially with distance and reaches $\approx 0.03$ at 1000~nm.
        Switching to the proximity-free weight vector $(0, 0.45, 0.40, 0.15)$
        for long-range pairs prevented the suppression of genuine long-range
        splits: 27 new ground-truth pairs that were previously filtered out by
        the \texttt{min\_composite\_score} floor (because proximity collapsed
        their score below the threshold) now entered the candidate pool,
        extending coverage by 1.8 percentage points (81.7\% $\to$ 83.5\%).
        Because these newly recovered pairs were harder cases (48\% accept
        rate vs 90\% for pre-existing pairs), aggregate recall moved slightly
        from 0.898 to 0.889 before subsequent tuning (Experiments 16--17)
        raised it again.

  \item \textbf{Experiment 17 (MaxDistance + SizeDiscrepancy):} PCA-based
        endpoint pairs can span up to $\sim$12{,}000~nm (the tip-to-tip distance
        across a large fragment's bounding box). The MaxDistanceRule at 5{,}000~nm
        was hard-rejecting these; raising it to 20{,}000~nm recovered all such
        pairs. Similarly, tiny stub fragments (100--200 voxels) attached to large
        trunks ($\sim$600{,}000 voxels) have cube-root radius ratios of 15--18,
        which required raising \texttt{max\_radius\_ratio} from 5.0 to 15.0.
        Together, these changes raised recall from 0.916 to 0.994.
\end{enumerate}

\paragraph{Negative experiments.}
Three approaches were tested and reverted:

\begin{itemize}
  \item \textbf{Experiment 13 (3000 nm radius):} Raising \texttt{max\_endpoint\_search\_nm}
        from 2000 to 3000 nm introduced additional long-range candidates that
        could not be distinguished from true splits, causing recall to regress
        from 0.898 to 0.865.

  \item \textbf{Experiment 18 (partner limit):} A post-acceptance filter that
        retained only the top-$N$ scoring partners per fragment was implemented
        with AND semantics (both fragments must rank each other in their top-$N$).
        This caused catastrophic recall collapse (Recall = 0.008): true merge
        pairs ranked at the median (position 46 of 63) by composite score within
        their fragment's partner list, with short-gap false positives scoring
        higher on proximity and ranking above them.

  \item \textbf{Experiment 21 (ML filter at 3000 nm):} A gradient-boosted
        classifier trained on 8 features (6 score components + per-fragment
        degree) at the 3000~nm radius setting achieved PR-AUC = 0.071 and could
        not recover recall $\geq 0.99$ at any threshold.
\end{itemize}

\section{Best Result: Experiment 24}

Experiment 24 introduces direction-weighted short-range scoring (raising alignment
weight from 0.30 to 0.45 and lowering proximity weight from 0.35 to 0.15),
disables the MinGapRule (\texttt{min\_gap\_nm} = 0), and raises
\texttt{max\_radius\_ratio} from 15.0 to 20.0. Together these changes recover all
37 MinGapRule false negatives and 2 SizeDiscrepancyRule false negatives from the
Experiment 19 configuration, reducing FN from 43 to 5.

\begin{table}[H]
\centering
\caption{XPRESS training volume: Experiment 24 final results.}
\label{tab:xpress_best}
\begin{tabular}{@{}lr@{}}
\toprule
Metric & Value \\
\midrule
Fragments & 11{,}208 \\
Candidates generated & 369{,}570 \\
Candidates accepted & 368{,}181 \\
Ground-truth pairs & 1{,}499 \\
Coverage & 78.7\% (1{,}180 / 1{,}499) \\
True Positives (TP) & 1{,}175 \\
False Positives (FP) & 366{,}933 \\
False Negatives (FN) & 5 \\
Recall & \textbf{0.9958} \\
Precision & $\approx$0.003 \\
\bottomrule
\end{tabular}
\end{table}

The 319 ground-truth pairs not covered by the candidate graph fall into two categories:
247 are structurally unreachable at the current 2000~nm search radius (172 pairs
with true gap $> 2000$~nm; 75 involving PCA-only large fragments whose tip-to-tip
distance does not span the gap), and the remaining 72 score below the acceptance
threshold under the revised weight vector. All 5 false negatives among covered
pairs fail the CurvatureRule (direction reversal exceeding 150$^\circ$, gap
66--987~nm), with composite scores of 0.41--0.64 that are above the acceptance
threshold but below CurvatureRule's hard reject.

\paragraph{The precision challenge.}
Precision $\approx$ 0.003 reflects a fundamental challenge of the XPRESS
domain: at the 2000~nm search radius, the pipeline generates 369{,}570
candidates, of which the vast majority are between different axons whose
local geometry
happens to be compatible. The pipeline's geometric features (alignment, continuity)
are informative for ruling out obviously incorrect pairs, but are insufficient
to discriminate same-axon from cross-axon pairs at long range.

\section{Machine Learning Filter (Experiment 20)}

To improve precision without sacrificing recall, a gradient-boosted classifier
(scikit-learn \texttt{GradientBoostingClassifier}) was trained on 8 features:
the 6 score components (proximity, alignment, continuity, size, composite,
gap\_distance) plus per-fragment \emph{degree} (number of accepted connections
for each fragment in the candidate list). The degree feature is the key
discriminator: fragments involved in genuine splits typically have 1--5 accepted
partners, while PCA-based large fragments generate 500--2{,}000 spurious accepted
connections.

\begin{table}[H]
\centering
\caption{ML filter (Experiment 20) at selected decision thresholds. PR-AUC = 0.098
vs.\ random baseline of 0.004.}
\label{tab:ml_filter}
\begin{tabular}{@{}rrrr@{}}
\toprule
Threshold & Recall & Precision & FP reduction \\
\midrule
0.0000 & 0.994 & 0.004 & 0\% (no filter) \\
0.0004 & 0.961 & 0.007 & 48.2\% \\
0.0050 & 0.852 & 0.037 & 91.9\% \\
0.0100 & 0.793 & 0.073 & 96.4\% \\
\bottomrule
\end{tabular}
\end{table}

The filter achieves a $10\times$ precision gain (to 0.037) at the cost of 14.2
percentage points of recall. No threshold achieves recall $\geq 0.99$ with
meaningful precision improvement, indicating that the classifier's decision
boundary is not well-aligned with the TP/FP boundary at high recall.

\section{False Negative Diagnosis (Experiment 22)}

A systematic diagnosis of false negatives under the current configuration
(with MinGapRule at 150~nm, corresponding to Experiment 19) is shown in
Table~\ref{tab:fn_diagnosis}.

\begin{table}[H]
\centering
\caption{FN attribution: which rule first rejects each false-negative candidate
in the current pipeline configuration (with MinGapRule 150 nm).}
\label{tab:fn_diagnosis}
\begin{tabular}{@{}llrl@{}}
\toprule
Rule & Count & Gap range & Notes \\
\midrule
MinGapRule (150 nm) & 37 & 33--148 nm & Genuine short-gap splits; recoverable \\
CurvatureRule & 4 & 209--987 nm & Hard cases; direction estimate unreliable \\
SizeDiscrepancyRule & 2 & 2189, 6683 nm & Stub vs.\ large trunk; fixable \\
\midrule
Total & 43 & & \\
\bottomrule
\end{tabular}
\end{table}

MinGapRule was the dominant false-negative cause (86\%) in this diagnostic
configuration. Experiment~24 resolved this by setting \texttt{min\_gap\_nm} = 0,
recovering all 37 MinGapRule FN candidates and 114 additional ground-truth pairs
(whose best candidates all had gap $< 150$~nm), and raising training recall
to 0.9958. The CurvatureRule FNs (4 cases here, 5 in Exp 24 after further
parameter changes) remain the only unresolved source of false negatives.

\section{Held-Out Validation Results (Experiments 23 and 24)}

To assess generalization, the pipeline was applied unchanged to the XPRESS
held-out validation volume ($699^3$ voxels, 203 ground-truth merge pairs, offset
(252, 252, 252) voxels within the full-volume coordinate system). Two
configurations were evaluated: Experiment 23 uses the Experiment 17 config
(MinGapRule active at 150~nm); Experiment 24 uses the final best config
(MinGapRule disabled, direction-weighted scoring).

\begin{table}[H]
\centering
\caption{Generalization comparison: training and held-out validation results for
two configurations.}
\label{tab:generalization}
\begin{tabular}{@{}lrrrr@{}}
\toprule
 & \multicolumn{2}{c}{Exp 23 config (MinGapRule on)} & \multicolumn{2}{c}{Exp 24 config (MinGapRule off)} \\
\cmidrule(lr){2-3}\cmidrule(lr){4-5}
Metric & Training & Validation & Training & Validation \\
\midrule
Ground-truth pairs & 1{,}499 & 203 & 1{,}499 & 203 \\
Coverage & 75.9\% & 72.4\% (147) & 78.7\% (1{,}180) & 81.3\% (165) \\
TP & 1{,}097 & 127 & 1{,}175 & 163 \\
FN & 43 & 20 & 5 & 2 \\
Recall & 0.966 & 0.885 & \textbf{0.9958} & \textbf{0.9879} \\
FN from CurvatureRule & 4 & 2 & 5 & 2 \\
\bottomrule
\end{tabular}
\end{table}

The FN composition confirms that removing the MinGapRule completely eliminates
the dominant false-negative source. Under the Exp 24 config, all remaining
false negatives on both volumes are caused by the CurvatureRule (5 on training,
gap 66--987~nm; 2 on validation, gap 382 and 693~nm), with no remaining
MinGapRule or SizeDiscrepancyRule rejections. The proportional similarity
between training and validation FN compositions (CurvatureRule-only in both)
confirms that the thresholds tuned on the training volume generalize without
overfitting.

\section{Visual Inspection}

Figures~\ref{fig:accepted} and~\ref{fig:rejected} show representative candidate
pairs from the XPRESS training volume using the three-panel layout described in
Section~\ref{sec:visual-validation}: neutral segmentation with skeleton overlay
(Panel~1), pipeline prediction with score breakdown (Panel~2), and ground-truth
ground-truth verdict (Panel~3).  Candidates are cropped to a 200$\times$200 voxel
window ($6{,}600\,\text{nm}$ on a side) centered at the midpoint between the two
fragment centroids, large enough to show both fragments and the gap between them.
Same-label pairs (two fragments of the same segment, which trivially score 1.0)
are excluded from sampling.

\begin{figure}[htbp]
  \centering
  \includegraphics[width=\textwidth,height=0.78\textheight,keepaspectratio]{candidate_showcase_accepted.png}
  \caption{High-confidence accepted candidates (highest composite scores among
  accepted, non-same-label connections) from the XPRESS training volume.
  \textbf{Left (Panel~1):} Neutral segmentation with distinct muted colors per
  segment; yellow dots/lines show ground-truth skeleton nodes near the Z-slice.
  \textbf{Center (Panel~2):} Fragment~A in red, Fragment~B in blue, with the
  proposed connection arrow (ACC badge) and all five component scores.
  \textbf{Right (Panel~3):} Ground-truth verdict (TP/FP/FN/TN badge); bright yellow
  skeleton highlights the neurons passing through both fragments, with a dashed
  line marking ground-truth-positive edge crossings. The FP outcomes for these candidates
  reflect the low precision of the pipeline at this operating point: with
  ${\sim}368{,}000$ accepted candidates and only ${\sim}1{,}494$ true positives,
  even geometrically compelling pairs are predominantly false accepts.  The
  visual report makes this trade-off immediately visible.}
  \label{fig:accepted}
\end{figure}

\begin{figure}[htbp]
  \centering
  \includegraphics[width=\textwidth,height=0.78\textheight,keepaspectratio]{candidate_showcase_rejected.png}
  \caption{Borderline rejected candidates (highest composite scores among all
  rejected, non-same-label connections).  These are the cases the pipeline
  declines to merge while scoring closest to acceptance.  All three shown are
  TN (correct rejections): the ground truth confirms no skeleton crossing between the
  two fragments, consistent with the pipeline's decision.  The skeleton overlay
  in Panel~1 makes the absence of a crossing evident: the ground-truth axon
  paths run through each fragment independently without bridging the gap.}
  \label{fig:rejected}
\end{figure}

The accepted panels (Figure~\ref{fig:accepted}) reveal an important characteristic
of the pipeline's operating regime: the composite score is optimised for recall,
not precision.  At the validation threshold used in Experiment~24, the pipeline
accepts ${\sim}368{,}000$ candidates to achieve ${\sim}0.996$ recall on the
$1{,}499$ ground-truth pairs, a design choice that requires a downstream filtering
step (e.g., the ML filter from Experiment~20) for deployment scenarios where
precision matters.  The borderline rejected panels (Figure~\ref{fig:rejected})
confirm that the rejection decisions are consistent with the geometric evidence:
cases where alignment and continuity scores are high but the ground truth confirms
no true split error exist.

\paragraph{Ground-truth TP inspection (end-to-end validation).}
Through Experiment~25 all evaluation was metric-driven (coverage, recall,
precision).  The ground-truth TP sampling category described in
Section~\ref{sec:visual-validation} closes this loop by filtering to the
intersection of pipeline-accepted and ground-truth-positive candidates and
asking, for each sampled row, whether the proposed merge is anatomically
real.  Running the inspection on the Experiment~24 training output (top
five ground-truth TPs by composite score, written to
\texttt{docs/visual\_report\_tp\_inspection.pdf}) confirmed that each
inspected TP is a genuine biological split: the ground-truth skeleton
visibly traverses both fragments, the fragments are morphologically
continuous across the gap, and the skeleton overlay in Panel~3 occupies
both label regions with a dashed ground-truth-positive indicator.  This is
the final non-automated layer of end-to-end validation and should be
re-run after any change that alters accepted-candidate composition.

\section{Gap-Stratified Precision Analysis}
\label{sec:gap-stratified}

The aggregate XPRESS precision of $\approx$0.003 on the Experiment~24 training
output is a single number that conceals large structural heterogeneity.  To
localize where the precision problem lives, we partitioned the 369{,}570
candidates into gap-distance bins and recomputed per-bin precision, recall,
and candidate volume using the same skeleton-edge-crossing ground-truth
construction described in Chapter~\ref{ch:methodology}.  This analysis is a
pure post-processing pass over the existing \texttt{connections.csv} export;
no pipeline re-run was required.  The analysis script is reproducible at
\texttt{scripts/gap\_stratified\_precision.py}.  Results are summarized in
Table~\ref{tab:gap_stratified} and Figure~\ref{fig:gap_stratified}.

\begin{table}[H]
\centering
\caption{Per-bin precision/recall on the Experiment~24 training volume.  ``TP''
and ``FP'' are candidate-level counts (each accepted candidate contributes
independently).  ``Candidates'' is total generated (accepted $+$ rejected)
within the bin.  Precision is computed from candidate counts; recall from
covered ground-truth pairs within the bin.  Note: the candidate-level TP total
of 1{,}248 differs from the 1{,}175 ``unique GT merge pairs accepted'' figure
reported in Table~\ref{tab:xpress_best} because a single ground-truth pair can
appear as multiple distinct accepted candidates (e.g., when several skeleton
edges of the same axon cross the fragment boundary).  Both accountings
describe the same underlying run.}
\label{tab:gap_stratified}
\small
\begin{tabular}{@{}lrrrrr@{}}
\toprule
Gap bin (nm) & Candidates & TP & FP & Precision & Recall \\
\midrule
$[0, 100)$       & 3{,}675    & 25    & 3{,}636     & 0.0068 & 0.96 \\
$[100, 200)$     & 3{,}352    & 44    & 3{,}283     & 0.0132 & 1.00 \\
$[200, 300)$     & 6{,}261    & 125   & 6{,}090     & 0.0201 & 0.99 \\
$[300, 500)$     & 20{,}418   & 594   & 19{,}745    & \textbf{0.0292} & 1.00 \\
$[500, 750)$     & 1{,}720    & 11    & 1{,}663     & 0.0066 & 1.00 \\
$[750, 1000)$    & 8{,}670    & 35    & 8{,}021     & 0.0043 & 0.95 \\
$[1000, 1500)$   & 52{,}062   & 117   & 51{,}936    & 0.0022 & 1.00 \\
$[1500, 2000)$   & \textbf{222{,}788} & 203 & \textbf{222{,}572} & 0.0009 & 1.00 \\
$[2000, 3000)$   & 1{,}902    & 32    & 1{,}856     & 0.0169 & 1.00 \\
$[3000, 20000)$  & 48{,}720   & 62    & 48{,}131    & 0.0013 & 1.00 \\
\midrule
Total            & 369{,}570  & 1{,}248 & 366{,}933 & 0.0034 & 0.9960 \\
\bottomrule
\end{tabular}
\end{table}

\begin{figure}[H]
\centering
\includegraphics[width=\textwidth]{gap_stratified_precision.png}
\caption{XPRESS Experiment~24 per-bin analysis. \textbf{Top:} in-bin precision
on a log scale, peaking at 0.029 in the $[300, 500)$~nm bin (about $10\times$
the overall figure) and collapsing to 0.0009 in the $[1500, 2000)$~nm bin.
\textbf{Bottom:} true-positive (green) and false-positive (red) candidate
counts per bin on a symlog scale. The $[1500, 2000)$~nm bin alone accounts for
60.7\% of all false positives (222{,}572 of 366{,}933) while contributing
only 16.3\% of true positives.}
\label{fig:gap_stratified}
\end{figure}

\paragraph{Findings.}
Three structural observations emerge:

\begin{enumerate}
  \item \textbf{Short-gap precision is an order of magnitude higher.}
        The $[100, 500)$~nm bins achieve precision 0.013--0.029, compared to
        the overall 0.003.  These gap sizes correspond to genuine interior
        skeleton-node splits where the two fragments are close enough that
        alignment and continuity scores are informative.

  \item \textbf{Long-range pairs dominate the FP budget.}
        The $[1000, 2000)$~nm range --- the regime where the distance-conditioned
        weight vector is active --- contains 74.4\% of all candidates
        (274{,}850 / 369{,}570) and 74.8\% of all false positives
        (274{,}508 / 366{,}933), but only 25.6\% of true positives
        (320 / 1{,}248).  The long-range pass pays a precision price that
        short-range scoring does not.

  \item \textbf{Precision rebounds at $[2000, 3000)$~nm}
        (0.017, second-highest of all bins) because these candidates are rare
        (1{,}902 total) and most arise from PCA-endpoint pairings on large
        fragments, where the bbox-tip estimate is only informative for
        genuine long axon splits.  Above 3{,}000~nm, spurious PCA pairings
        dominate and precision drops again.
\end{enumerate}

\paragraph{Generalization to validation volume.}
Running the same post-processing on the Experiment~24 held-out validation
output reproduces the same shape: the $[300, 500)$~nm bin is the precision
peak (0.0025) and the $[1500, 2000)$~nm bin alone contains 65.5\% of all
false positives (300{,}763 / 458{,}941).  The absolute precision numbers
are lower on validation because there are fewer ground-truth pairs per
candidate (203 vs.\ 1{,}499), but the relative ordering across bins is
preserved.  This confirms that the precision-by-gap structure is a
property of the XPRESS domain (white matter axon density at long range),
not an artifact of training-volume tuning.

\paragraph{Implications for future precision work.}
Three directions follow directly from the per-bin decomposition:

\begin{itemize}
  \item A gap-dependent accept threshold --- looser at $[100, 500)$~nm, much
        stricter at $[1000, 2000)$~nm --- would retain most TPs while cutting
        most FPs.  The $[300, 500)$~nm bin retains 594 TPs at precision 0.029
        with a uniform threshold; moving the $[1500, 2000)$~nm threshold up
        to discard its 222{,}572 FPs would cost at most 203 TPs.
  \item The machine-learned filter of Experiment~20 should be retrained
        stratified by gap, since the feature distribution and operating
        point differ sharply between bins.
  \item Reducing \texttt{max\_endpoint\_search\_nm} from 2000 to 1000~nm is
        worth re-testing.  Experiment~13 (3000~nm) was rejected for recall
        regression, but the reverse direction --- tightening to 1000~nm ---
        would discard the long-range bins entirely, lose the 117 TPs at
        $[1000, 1500)$~nm and 203 at $[1500, 2000)$~nm (25.6\% of total TPs),
        and save $\sim$270{,}000 FPs.  This would yield recall near 0.73
        at precision $\approx$0.015, a substantially different operating
        point that may be preferable when precision matters more than
        absolute recall.
\end{itemize}

\section{Neighborhood-Density Re-ranker}
\label{sec:density-reranker}

The per-bin analysis (Section~\ref{sec:gap-stratified}) localizes the
precision problem to the long-range regime but does not explain why one
long-range candidate is a true split and another is not.  A natural
hypothesis --- motivated by the direction-weighted scoring story of
Experiment~24 --- is that a genuine same-axon split pair is
\emph{locally isolated}: its axon is the primary structure in the
neighborhood.  A cross-axon false positive sits in a crowd of other axons
that makes many geometrically-plausible alternatives to the pipeline's
chosen partner.  A simple count of other fragments in the neighborhood
of the candidate should therefore separate TPs from FPs even after the
composite score has already accepted both.

This section tests that hypothesis as a post-hoc re-ranker on the
Experiment~24 accepted-candidate set.  The analysis script
(\texttt{scripts/neighborhood\_density\_reranker.py}) is pure
post-processing: no pipeline code or pipeline configuration is modified.

\paragraph{Feature definition.}
For each accepted candidate with fragments $a$ and $b$:
\begin{align*}
\text{mid}    &= (\text{centroid}_a + \text{centroid}_b) / 2 \\
\text{density\_mid}(R) &= \big|\{\,c : c \neq a, b \text{ and } \|\text{centroid}_c - \text{mid}\| \le R\,\}\big| \\
\text{density\_max}(R) &= \max\left(|\mathcal{N}_R(a)|,\, |\mathcal{N}_R(b)|\right) - 2
\end{align*}
where $\mathcal{N}_R(f)$ is the set of fragment centroids within radius $R$
of fragment $f$'s centroid.  We evaluated $R \in \{500, 1000, 1500, 2000, 3000\}$~nm;
$R = 1500$~nm gave the best precision-at-recall scores and is used in the
figures below.

\paragraph{Distribution of the feature.}
Figure~\ref{fig:density-distribution} shows the per-candidate density
distributions stratified by ground-truth status.  The TP distribution is
visibly shifted toward lower density: for \texttt{density\_mid\_1500}, TP
mean is 9.6 vs FP mean 13.4; for \texttt{density\_max\_1500}, TP mean is
12.1 vs FP mean 15.6.  The two distributions overlap substantially but
the means are distinct, confirming the qualitative hypothesis.

\begin{figure}[H]
\centering
\includegraphics[width=\textwidth]{density_distribution.png}
\caption{Neighborhood-density feature distributions for the 368{,}181
Experiment~24 accepted candidates, stratified by ground-truth status
(TP = accepted $\land$ GT-positive, FP = accepted $\land$ $\neg$GT).
\textbf{Left:} \texttt{density\_mid\_1500} (count of fragment centroids
within 1{,}500~nm of the candidate midpoint). \textbf{Right:}
\texttt{density\_max\_1500} (max of the per-fragment counts). The TP
distribution is shifted toward lower density in both cases.}
\label{fig:density-distribution}
\end{figure}

\paragraph{Re-ranker performance.}
Table~\ref{tab:density-reranker} reports precision at fixed recall
targets for three ranking strategies: the baseline composite score
(descending), the density features alone (ascending), and a combined
score that adds $\alpha \log(1 + \text{density\_mid})$ as a penalty
term.  The composite-score baseline within the already-accepted set is
remarkably uninformative --- precision varies by less than 3\% across
the 0.85--0.99 recall range --- reflecting that all accepted candidates
occupy a narrow composite-score band by construction.  The density
features provide a genuine if modest improvement at mid-recall,
reaching a 49\% relative gain at recall $0.85$.

\begin{table}[H]
\centering
\caption{Precision at fixed recall targets on Experiment~24 accepted
candidates (training volume).  Numbers in parentheses show the validation
volume.  ``Rel.'' is the relative improvement of
\texttt{density\_max\_1500} over the composite-score baseline on training.}
\label{tab:density-reranker}
\small
\begin{tabular}{@{}lrrrr@{}}
\toprule
Recall & Composite baseline & density\_mid\_1500 & density\_max\_1500 & Rel. \\
\midrule
0.99 & 0.0034 \,(0.00038) & 0.0036 \,(0.00040) & 0.0035 \,(0.00040) & $+$4\% \\
0.95 & 0.0034 \,(0.00038) & 0.0039 \,(0.00042) & 0.0039 \,(0.00045) & $+$15\% \\
0.90 & 0.0034 \,(0.00038) & 0.0043 \,(0.00044) & 0.0045 \,(0.00054) & $+$33\% \\
0.85 & 0.0033 \,(0.00037) & 0.0048 \,(0.00048) & 0.0049 \,(0.00056) & $+$49\% \\
\bottomrule
\end{tabular}
\end{table}

\begin{figure}[H]
\centering
\includegraphics[width=0.92\textwidth]{density_pr_curve.png}
\caption{Precision-recall curves on the 368{,}181 Experiment~24 accepted
candidates for four ranking strategies: composite score (blue, baseline),
density\_mid\_1500 (red), density\_max\_1500 (green), and a combined
score composite $+ 0.1 \log(1 + \text{density\_mid\_1500})$ (purple).
Density features dominate the baseline across the 0.2--0.95 recall
range.  Note the log-scale on precision.  All curves converge at
recall $=1$, where every accepted candidate is kept and precision is
the aggregate 0.0034.}
\label{fig:density-pr}
\end{figure}

\paragraph{Findings.}

\begin{enumerate}
  \item \textbf{The hypothesis is empirically confirmed.}  Accepted-candidate
        density distributions differ between TP and FP populations
        (TP mean $\approx 25\%$ lower than FP mean at $R=1500$~nm), and a
        density-only re-ranker improves precision over composite-score
        re-ranking by up to 49\% at recall 0.85.

  \item \textbf{The effect decays sharply at high recall.}  At recall 0.99
        the improvement is only 4--6\% in both training and validation,
        because the $\sim$1\% of TPs that are hardest to recover have density
        indistinguishable from the FP background.  Density alone does not
        break the geometric-features precision ceiling at the recall target
        where the pipeline currently operates.

  \item \textbf{Composite score is not a useful re-ranker for already-accepted
        candidates.}  Its PR curve within the accepted set is nearly flat at
        precision $\approx 0.003$ across the full recall range
        (Figure~\ref{fig:density-pr}).  All accepted candidates have, by
        construction, composite scores above the acceptance threshold, and
        within that band score magnitude carries little signal about
        ground-truth status.

  \item \textbf{Density is complementary, not redundant, with the degree
        feature of Experiment~20.}  The ML filter of Experiment~20 achieved
        precision 0.037 at recall 0.85 using per-fragment degree.  Density
        at recall 0.85 achieves only 0.005 on its own --- an order of
        magnitude weaker --- but the two features encode different signals
        (how many partners has this fragment been paired with vs.\ how
        crowded is the local neighborhood).  The combined composite + density
        score (purple curve) does \emph{not} beat density alone in our
        experiments, because composite is a pure noise term on this slice;
        a learned combination with the degree feature is expected to do
        better and is left for future work.

  \item \textbf{Validation reproduces the pattern.}  On the held-out
        validation volume, the same relative improvement ordering holds
        (+6\% at recall 0.99, +50\% at recall 0.85), confirming that
        density is a real signal in this domain rather than a
        training-set artifact.
\end{enumerate}

\paragraph{Implication.}
The neighborhood-density feature is a worthwhile addition to any future
learned post-filter but is insufficient as a standalone precision
recovery strategy at recall $\geq 0.99$.  This is consistent with the
overall narrative of the XPRESS precision challenge: no single
geometric signal separates same-axon from cross-axon pairs at the
density of myelinated white matter, and meaningful precision gains at
high recall will require either multi-feature ML (retrained on features
including density) or access to raw image data.

\section{Per-Bin Threshold Precision Ceiling}
\label{sec:per-bin-threshold}

Sections~\ref{sec:gap-stratified} and~\ref{sec:density-reranker}
characterized the precision problem and attempted a single-feature
remedy.  A natural follow-up question is: \emph{what is the best
precision any pipeline variant can achieve if its only degree of
freedom is a per-gap-bin composite threshold?}  This sets an empirical
ceiling on all gap-aware threshold-tuning strategies --- including
the gap-dependent accept thresholds suggested by Section~\ref{sec:gap-stratified}
and any future ML filter whose inputs are limited to the current
geometric scores plus gap bin.  Quantitative results are summarized in
Table~\ref{tab:per-bin-threshold} and the full PR frontier is shown
in Figure~\ref{fig:per-bin-threshold}.

\paragraph{Method.}
For each accepted candidate we know the gap bin, the composite score,
and the ground-truth label.  The per-bin ceiling is constructed
by sorting candidates within each gap bin by composite score
(descending) and then merging the per-bin streams greedily: at each
step, whichever bin has a ground-truth positive at its next pointer
contributes that candidate for free; when no such bin exists, the bin
with the highest remaining TP-density (remaining TPs / remaining
candidates) contributes its next candidate.  This produces a global
PR curve that is the upper envelope across all per-bin
threshold-assignment strategies.  The construction uses ground-truth
labels to determine the merge order, so it represents an upper bound;
any learnable per-bin policy will sit on or below this curve.

\begin{table}[H]
\centering
\caption{Precision at fixed recall targets for the global composite
threshold (Experiment~24 baseline) vs the per-bin ground-truth ceiling.
Numbers in parentheses are the held-out validation volume.
``Rel.~gain'' is the relative precision improvement of per-bin vs
global on the training volume.}
\label{tab:per-bin-threshold}
\small
\begin{tabular}{@{}lrrr@{}}
\toprule
Recall & Global baseline & Per-bin ceiling & Rel.~gain \\
\midrule
0.99 & 0.0034 \,(0.00038) & 0.0034 \,(0.00041) & $+$1.3\% \\
0.95 & 0.0034 \,(0.00038) & 0.0040 \,(0.00055) & $+$16.4\% \\
0.90 & 0.0034 \,(0.00038) & 0.0051 \,(0.00087) & $+$50.6\% \\
0.85 & 0.0033 \,(0.00037) & 0.0068 \,(0.00104) & $+$104\% \\
0.80 & 0.0033 \,(0.00037) & 0.0085 \,(0.00143) & $+$161\% \\
0.70 & 0.0032 \,(0.00035) & \textbf{0.0169} \,(\textbf{0.00303}) & $+$421\% \\
\bottomrule
\end{tabular}
\end{table}

\begin{figure}[H]
\centering
\includegraphics[width=0.92\textwidth]{per_bin_threshold_pr.png}
\caption{Precision-recall frontier for the two strategies on the
Experiment~24 accepted-candidate set.  \textbf{Blue:} the global
composite threshold (the baseline behavior of Experiment~24 after the
fact) is nearly flat at precision $\approx 0.003$ across the full
recall range.  \textbf{Red:} the per-bin ground-truth ceiling reaches
$\sim$4\% precision at low recall and holds above 1\% up to recall
$\approx 0.80$ before collapsing to meet the baseline at recall 1.
The two curves converge at recall $1$ where every accepted candidate
is kept.  Log-scale precision.}
\label{fig:per-bin-threshold}
\end{figure}

\paragraph{Findings.}

\begin{enumerate}
  \item \textbf{The recall-0.99 ceiling is nearly flat.}  At recall 0.99
        the per-bin ceiling gains only 1.3\% over the global baseline
        (9\% on validation).  The hardest-to-recover 1\% of TPs are
        scattered across the same bins that hold the bulk of the FPs, so
        no per-bin threshold allocation can cut FPs meaningfully without
        sacrificing those TPs.  This explains why Experiment~20's ML
        filter could not maintain recall $\geq 0.99$ with meaningful
        precision improvement: the ceiling is intrinsic to the
        geometric feature space plus gap bin, not an ML-architecture
        shortcoming.

  \item \textbf{The recall-$\leq 0.90$ ceiling is substantially higher.}
        At recall 0.90, per-bin thresholding can reach 0.005 precision
        ($+51\%$); at recall 0.85, 0.007 ($+104\%$); at recall 0.70,
        0.017 ($+421\%$).  These are the operating points at which
        per-bin threshold tuning becomes a meaningful deployment lever.

  \item \textbf{Validation reproduces the structure.}  On the held-out
        volume the absolute precision numbers are lower, but the
        relative gain ordering is preserved: +9\% at recall 0.99 rising
        to +773\% at recall 0.70.

  \item \textbf{Comparing to alternative precision strategies.}
        Experiment~20's ML filter reached precision 0.037 at recall
        0.85.  The per-bin ceiling at recall 0.85 is 0.007 --- 5x
        lower --- meaning Experiment~20's degree-feature signal gave
        access to information beyond what any gap-aware threshold
        strategy can capture.  Conversely, the density re-ranker of
        Section~\ref{sec:density-reranker} peaked at 0.005 at recall
        0.85, which is close to the per-bin ceiling: density
        plus gap captures nearly all the gap-bin-aware precision
        signal, and further gains require orthogonal features such
        as per-fragment degree or raw-image evidence.
\end{enumerate}

\paragraph{Implication.}
Per-bin threshold tuning is a concrete, low-risk deployment lever at
recall $\leq 0.90$ but delivers essentially nothing at recall 0.99.
Any future work targeting high-recall precision must combine
orthogonal feature classes (per-fragment degree, neighborhood density,
skeleton-morphology features, or raw image evidence) rather than
merely redistributing the existing composite score.  This result
constrains the design of the recommended multi-feature gap-stratified
ML filter (Chapter~\ref{ch:futurework}): a useful model must move
decisions outside the convex hull set by composite-plus-gap, not simply
calibrate within it.

\section{Test Coverage}

The pipeline is validated by 428 automated tests across 20 test modules,
covering all pipeline stages, configuration loading, export formats, and edge
cases. All tests pass on Python 3.10, 3.11, and 3.12. Continuous integration
via GitHub Actions runs tests, black formatting checks, and mypy type checks on
every push.

\begin{table}[H]
\centering
\caption{Test suite summary.}
\label{tab:tests}
\begin{tabular}{@{}lr@{}}
\toprule
Module & Tests \\
\midrule
I/O readers & 22 \\
Fragment extraction & 35 \\
Graph construction & 18 \\
Candidate generation \& scoring & 40 \\
Validation rules & 85 \\
Assembly & 35 \\
Export formats & 50 \\
Ground truth evaluation & 30 \\
ML filter & 28 \\
Visualization & 12 \\
Configuration & 28 \\
End-to-end pipeline & 12 \\
Other (CLI, logging, spatial math) & 33 \\
\midrule
\textbf{Total} & \textbf{428} \\
\bottomrule
\end{tabular}
\end{table}
